% Figure from MATH 452 notes, section 5. Compile with the notes preamble.
\begin{tikzpicture}[scale=1.3,
  mid/.style={postaction={decorate}, decoration={markings, mark=at position 0.55 with {\arrow{stealth}}}}]
\draw[-stealth, thin] (-0.2,0) -- (4.9,0) node[right, font=\scriptsize] {$x$};
\draw[-stealth, thin] (0,-0.2) -- (0,2.95) node[above, font=\scriptsize] {$y$};
% rectangle corners
\coordinate (A) at (1.5,0.7);
\coordinate (B) at (3.7,0.7);
\coordinate (C) at (3.7,2.2);
\coordinate (D) at (1.5,2.2);
\fill[figB!7] (A) rectangle (C);
% phi = difference in x (horizontal edges), psi = difference in y (vertical edges)
\draw[figG, thick, mid] (A) -- (B) node[midway, below, font=\scriptsize, inner sep=3pt] {$\varphi(y_0)$};
\draw[figG, thick, mid] (D) -- (C) node[midway, above, font=\scriptsize, inner sep=3pt] {$\varphi(y_0+k)$};
\draw[figO, thick, mid] (A) -- (D) node[midway, left, font=\scriptsize, inner sep=3pt] {$\psi(x_0)$};
\draw[figO, thick, mid] (B) -- (C) node[midway, right, font=\scriptsize, inner sep=3pt] {$\psi(x_0+h)$};
% corners
\filldraw (A) circle (1.2pt) node[below left, font=\scriptsize, inner sep=1.5pt] {$(x_0,y_0)$};
\filldraw (B) circle (1.2pt) node[below right, font=\scriptsize, inner sep=1.5pt] {$(x_0+h,\,y_0)$};
\filldraw (C) circle (1.2pt) node[above right, font=\scriptsize, inner sep=1.5pt] {$(x_0+h,\,y_0+k)$};
\filldraw (D) circle (1.2pt) node[above left, font=\scriptsize, inner sep=1.5pt] {$(x_0,\,y_0+k)$};
% signs of the mixed second difference
\node[black!60, font=\scriptsize] at ($(A)+(0.17,0.17)$) {$+$};
\node[black!60, font=\scriptsize] at ($(B)+(-0.17,0.17)$) {$-$};
\node[black!60, font=\scriptsize] at ($(C)+(-0.17,-0.17)$) {$+$};
\node[black!60, font=\scriptsize] at ($(D)+(0.17,-0.17)$) {$-$};
% the two mean-value points where I(h,k) is attained
\filldraw[figR] (2.25,1.15) circle (1.3pt) node[right, font=\scriptsize, inner sep=2pt] {$f_{xy}$};
\filldraw[figR] (3.05,1.75) circle (1.3pt) node[left, font=\scriptsize, inner sep=2pt] {$f_{yx}$};
\end{tikzpicture}
